% Options for packages loaded elsewhere
% Options for packages loaded elsewhere
\PassOptionsToPackage{unicode}{hyperref}
\PassOptionsToPackage{hyphens}{url}
\PassOptionsToPackage{dvipsnames,svgnames,x11names}{xcolor}
%
\documentclass[
  letterpaper,
  DIV=11,
  numbers=noendperiod]{scrartcl}
\usepackage{xcolor}
\usepackage{amsmath,amssymb}
\setcounter{secnumdepth}{-\maxdimen} % remove section numbering
\usepackage{iftex}
\ifPDFTeX
  \usepackage[T1]{fontenc}
  \usepackage[utf8]{inputenc}
  \usepackage{textcomp} % provide euro and other symbols
\else % if luatex or xetex
  \usepackage{unicode-math} % this also loads fontspec
  \defaultfontfeatures{Scale=MatchLowercase}
  \defaultfontfeatures[\rmfamily]{Ligatures=TeX,Scale=1}
\fi
\usepackage{lmodern}
\ifPDFTeX\else
  % xetex/luatex font selection
\fi
% Use upquote if available, for straight quotes in verbatim environments
\IfFileExists{upquote.sty}{\usepackage{upquote}}{}
\IfFileExists{microtype.sty}{% use microtype if available
  \usepackage[]{microtype}
  \UseMicrotypeSet[protrusion]{basicmath} % disable protrusion for tt fonts
}{}
\makeatletter
\@ifundefined{KOMAClassName}{% if non-KOMA class
  \IfFileExists{parskip.sty}{%
    \usepackage{parskip}
  }{% else
    \setlength{\parindent}{0pt}
    \setlength{\parskip}{6pt plus 2pt minus 1pt}}
}{% if KOMA class
  \KOMAoptions{parskip=half}}
\makeatother
% Make \paragraph and \subparagraph free-standing
\makeatletter
\ifx\paragraph\undefined\else
  \let\oldparagraph\paragraph
  \renewcommand{\paragraph}{
    \@ifstar
      \xxxParagraphStar
      \xxxParagraphNoStar
  }
  \newcommand{\xxxParagraphStar}[1]{\oldparagraph*{#1}\mbox{}}
  \newcommand{\xxxParagraphNoStar}[1]{\oldparagraph{#1}\mbox{}}
\fi
\ifx\subparagraph\undefined\else
  \let\oldsubparagraph\subparagraph
  \renewcommand{\subparagraph}{
    \@ifstar
      \xxxSubParagraphStar
      \xxxSubParagraphNoStar
  }
  \newcommand{\xxxSubParagraphStar}[1]{\oldsubparagraph*{#1}\mbox{}}
  \newcommand{\xxxSubParagraphNoStar}[1]{\oldsubparagraph{#1}\mbox{}}
\fi
\makeatother


\usepackage{longtable,booktabs,array}
\usepackage{calc} % for calculating minipage widths
% Correct order of tables after \paragraph or \subparagraph
\usepackage{etoolbox}
\makeatletter
\patchcmd\longtable{\par}{\if@noskipsec\mbox{}\fi\par}{}{}
\makeatother
% Allow footnotes in longtable head/foot
\IfFileExists{footnotehyper.sty}{\usepackage{footnotehyper}}{\usepackage{footnote}}
\makesavenoteenv{longtable}
\usepackage{graphicx}
\makeatletter
\newsavebox\pandoc@box
\newcommand*\pandocbounded[1]{% scales image to fit in text height/width
  \sbox\pandoc@box{#1}%
  \Gscale@div\@tempa{\textheight}{\dimexpr\ht\pandoc@box+\dp\pandoc@box\relax}%
  \Gscale@div\@tempb{\linewidth}{\wd\pandoc@box}%
  \ifdim\@tempb\p@<\@tempa\p@\let\@tempa\@tempb\fi% select the smaller of both
  \ifdim\@tempa\p@<\p@\scalebox{\@tempa}{\usebox\pandoc@box}%
  \else\usebox{\pandoc@box}%
  \fi%
}
% Set default figure placement to htbp
\def\fps@figure{htbp}
\makeatother
\usepackage{svg}


% definitions for citeproc citations
\NewDocumentCommand\citeproctext{}{}
\NewDocumentCommand\citeproc{mm}{%
  \begingroup\def\citeproctext{#2}\cite{#1}\endgroup}
\makeatletter
 % allow citations to break across lines
 \let\@cite@ofmt\@firstofone
 % avoid brackets around text for \cite:
 \def\@biblabel#1{}
 \def\@cite#1#2{{#1\if@tempswa , #2\fi}}
\makeatother
\newlength{\cslhangindent}
\setlength{\cslhangindent}{1.5em}
\newlength{\csllabelwidth}
\setlength{\csllabelwidth}{3em}
\newenvironment{CSLReferences}[2] % #1 hanging-indent, #2 entry-spacing
 {\begin{list}{}{%
  \setlength{\itemindent}{0pt}
  \setlength{\leftmargin}{0pt}
  \setlength{\parsep}{0pt}
  % turn on hanging indent if param 1 is 1
  \ifodd #1
   \setlength{\leftmargin}{\cslhangindent}
   \setlength{\itemindent}{-1\cslhangindent}
  \fi
  % set entry spacing
  \setlength{\itemsep}{#2\baselineskip}}}
 {\end{list}}
\usepackage{calc}
\newcommand{\CSLBlock}[1]{\hfill\break\parbox[t]{\linewidth}{\strut\ignorespaces#1\strut}}
\newcommand{\CSLLeftMargin}[1]{\parbox[t]{\csllabelwidth}{\strut#1\strut}}
\newcommand{\CSLRightInline}[1]{\parbox[t]{\linewidth - \csllabelwidth}{\strut#1\strut}}
\newcommand{\CSLIndent}[1]{\hspace{\cslhangindent}#1}



\setlength{\emergencystretch}{3em} % prevent overfull lines

\providecommand{\tightlist}{%
  \setlength{\itemsep}{0pt}\setlength{\parskip}{0pt}}



 


\KOMAoption{captions}{tableheading}
\makeatletter
\@ifpackageloaded{caption}{}{\usepackage{caption}}
\AtBeginDocument{%
\ifdefined\contentsname
  \renewcommand*\contentsname{Table of contents}
\else
  \newcommand\contentsname{Table of contents}
\fi
\ifdefined\listfigurename
  \renewcommand*\listfigurename{List of Figures}
\else
  \newcommand\listfigurename{List of Figures}
\fi
\ifdefined\listtablename
  \renewcommand*\listtablename{List of Tables}
\else
  \newcommand\listtablename{List of Tables}
\fi
\ifdefined\figurename
  \renewcommand*\figurename{Figure}
\else
  \newcommand\figurename{Figure}
\fi
\ifdefined\tablename
  \renewcommand*\tablename{Table}
\else
  \newcommand\tablename{Table}
\fi
}
\@ifpackageloaded{float}{}{\usepackage{float}}
\floatstyle{ruled}
\@ifundefined{c@chapter}{\newfloat{codelisting}{h}{lop}}{\newfloat{codelisting}{h}{lop}[chapter]}
\floatname{codelisting}{Listing}
\newcommand*\listoflistings{\listof{codelisting}{List of Listings}}
\makeatother
\makeatletter
\makeatother
\makeatletter
\@ifpackageloaded{caption}{}{\usepackage{caption}}
\@ifpackageloaded{subcaption}{}{\usepackage{subcaption}}
\makeatother
\usepackage{bookmark}
\IfFileExists{xurl.sty}{\usepackage{xurl}}{} % add URL line breaks if available
\urlstyle{same}
\hypersetup{
  pdftitle={Origin Ambiguity},
  pdfauthor={Sabou Rani Stocker},
  colorlinks=true,
  linkcolor={blue},
  filecolor={Maroon},
  citecolor={Blue},
  urlcolor={Blue},
  pdfcreator={LaTeX via pandoc}}


\title{Origin Ambiguity}
\author{Sabou Rani Stocker}
\date{2025-11-21}
\begin{document}
\maketitle


\section{Abstract}\label{abstract}

As firms increasingly use generative AI to produce product shoots, video
ads, and news-like reels, AI-generated outputs reach near-human quality,
making it difficult for viewers to distinguish whether content was
created by humans or AI. We take a new perspective on this phenomenon
and propose \emph{origin ambiguity} as a central challenge in such
contemporary marketing communication. In the absence of clear
regulations, consumers increasingly experience ambiguity related to
content origin and pose the question ``is this AI or real?''. We ask:
what changes when people no longer know if what they are looking at has
been generated by a human or by AI? To the best of our knowledge, this
question remains unanswered. Existing work suggests that ambiguity is a
relevant phenomenon with potentially detrimental effects. Specifically,
ambiguity might shift attention away from \emph{content deliberation}
toward deliberation about \emph{origin}, if consumers are motivated
enough to engage in deeper processing (e.g., under high involvement),
with implications for marketing persuasiveness, informed
decision-making, and affective responses (amongst others).

\section{Related theory and
phenomenon-in-context}\label{related-theory-and-phenomenon-in-context}

\subsection{AI and Human Origin}\label{ai-and-human-origin}

Much research has studied the consequences of explicit human vs.~AI
origin under the phenomena of algorithm aversion and appreciation (e.g.,
Bigman (2022)). Here, consumers know and/or are informed about the
origin of a recommendation, post, or product, and seem to prefer human
over AI or vice versa conditional on various factors. Other research
demonstrates that people are generally bad at identifying the origin of
content (AI or human) (e.g.,~Vukojičić, Krstić, and Veinović (2023);
Porter and Machery (2024); Milička et al. (2025)).

This latter research points at a human-AI generation frontier that is
increasingly blurred. As a case in point, multiple forums exist around
the question ``Is this AI or real/human''
(\href{https://www.reddit.com/r/RealOrAI/}{r/RealOrAI}), and the comment
sections on social media pages are filled with consumers eager to know
if what they are looking at is ``real'' or AI-generated. Importantly,
there is no legally mandated disclosure of AI involvement in content
generation, resulting in many situations where the origin remains
ambiguous.

\subsection{Ambiguity in Marketing}\label{ambiguity-in-marketing}

Ambiguity in marketing (e.g.~surprise boxes) has been shown to affect
attention and raise interest, but also result in aversion, especially in
high-stakes situations (Kovacheva and Nikolova 2024). For example, when
product quality is ambiguous, consumers rely more on ads (Hoch and Ha
1986), and on top-down approaches (i.e.~expectations, theories, Yi
(1993)); and ambiguity aversion correlates with a preference for
established products (Muthukrishnan, Wathieu, and Xu 2009; Hoch and Ha
1986). Furthermore, research on a distinct, but related phenomenon --
the uncanny valley (Mori, MacDorman, and Kageki 2012) -- gives
preliminary evidence that ambiguous stimuli tax cognitive resources
(Shin, Kim, and Biocca 2019; Tao et al. 2022; Howard 2017).

One important moderator for the occurrence of ambiguity effects in
general are the stakes of a situation and the resulting motivation of
consumers to resolve ambiguity. Marketing tactics strategically using
obfuscated materials (i.e.~cropping objects, incomplete typefaces,
\ldots) show that consumers attempt to resolve these introduced
ambiguities only when they are motivated to do so Hagtvedt (2011). We
expect that comparable effects occur under origin ambiguity.

\section{Relevance}\label{relevance}

\subsection{Theoretical}\label{theoretical}

This phenomenon integrates several streams of research that to the best
of our knowledge have been treated separately so far: research on
ambiguity in marketing, on cognitive load theory, and on consumer
exposure and reactions to AI-generated content. Extending prior work on
AI where origin is known or the variable in question, we study cases
where humans face content that is neither declared AI or human or easily
determined as one or the other.

\subsection{Practical}\label{practical}

\emph{Origin ambiguity} holds implications for both marketers and
consumers: For marketers, potential backfiring of AI generated ads as
well as the resulting trade-offs with efficiency and effectiveness are
relevant. For consumers, cognitive resource reallocation might come at
the cost of informed decision-making.

\section{Research Plan}\label{research-plan}

Based on the existing literature, we suggest that when people are no
longer sure if something has been generated by a human or by AI, they
try to resolve this ambiguity. One possible strategy here is more
elaboration on the origin and less deliberation on the content resulting
in worsened information retention. Additionally or alternatively, we
suggest that the ambiguity leads to a negative affective response.
Evidenced by previous research, we suggest that people only attempt to
resolve ambiguity when they have sufficient motivation to do so, in
other words when the stakes are high enough.

\begin{figure}

\centering{

\pandocbounded{\includesvg[keepaspectratio]{conceptual_model.svg}}

}

\caption{\label{fig-conceptual-model}Conceptual Model}

\end{figure}%

\section{Open Questions}\label{open-questions}

\begin{itemize}
\tightlist
\item
  Is this relevant and theoretically interesting enough in your eyes?
\item
  We mainly believe the IV (origin ambiguity) is interesting to study.
  We're less confident that what we propose downstream in terms of the
  ambiguity resolve strategy (cognitive resources allocated to origin vs
  content) is the best option available. Maybe you have other ideas what
  to look at?
\item
  Methodically, one crucial question pertains to the operationalization
  of the IV. Put simply, people may simply not think about origin
  without a prompt to do so. Specifically, we were wondering how we can
  test an origin ambiguity IV \emph{as it is naturally occuring} as
  opposed to pointing participants to the ambiguity of some content
  (which is then a confound).
\item
  A second question to be discussed is whether or not to include the
  comparison of low origin-ambiguity materials that are clearly AI
  generated, or if it makes more sense to only compare against clearly
  human-generated material. (we'd probably need to do so at some point)
\end{itemize}

\phantomsection\label{refs}
\begin{CSLReferences}{1}{0}
\bibitem[\citeproctext]{ref-bigmanAlgorithmicDiscriminationCauses2022}
Bigman, Yochanan E. 2022. {``Algorithmic {Discrimination Causes Less
Moral Outrage Than Human Discrimination}.''}

\bibitem[\citeproctext]{ref-hagtvedtImpactIncompleteTypeface2011}
Hagtvedt, Henrik. 2011. {``The {Impact} of {Incomplete Typeface Logos}
on {Perceptions} of the {Firm}.''} \emph{Journal of Marketing} 75 (4):
86--93. \url{https://doi.org/10.1509/jmkg.75.4.86}.

\bibitem[\citeproctext]{ref-hochConsumerLearningAdvertising1986}
Hoch, Stephen J., and Young-Won Ha. 1986. {``Consumer {Learning}:
{Advertising} and the {Ambiguity} of {Product Experience}.''}
\emph{Journal of Consumer Research} 13 (2): 221.
\url{https://doi.org/10.1086/209062}.

\bibitem[\citeproctext]{ref-howardInvestigatingSimulationElements2017}
Howard, Matt C. 2017. {``Investigating the Simulation Elements of
Environment and Control: {Extending} the {Uncanny Valley Theory} to
Simulations.''} \emph{Computers \& Education} 109 (June): 216--32.
\url{https://doi.org/10.1016/j.compedu.2017.03.005}.

\bibitem[\citeproctext]{ref-joharConsumerInvolvementDeception1995}
Johar, Gita Venkataramani. 1995. {``Consumer {Involvement} and
{Deception} from {Implied Advertising Claims}.''} \emph{JOURNAL OF
MARKETING RESEARCH}.

\bibitem[\citeproctext]{ref-kovachevaUncertaintyMarketingTactics2024}
Kovacheva, Aleksandra, and Hristina Nikolova. 2024. {``Uncertainty
Marketing Tactics: {An} Overview and a Unifying Framework.''}
\emph{Journal of the Academy of Marketing Science} 52 (1): 1--22.
\url{https://doi.org/10.1007/s11747-023-00941-7}.

\bibitem[\citeproctext]{ref-milickaHumansCanLearn2025}
Milička, Jiří, Anna Marklová, Ondřej Drobil, and Eva Pospíšilová. 2025.
{``Humans Can Learn to Detect {AI-generated} Texts, or at Least Learn
When They Can't.''} \emph{PLOS One} 20 (10): e0333007.
\url{https://doi.org/10.1371/journal.pone.0333007}.

\bibitem[\citeproctext]{ref-moriUncannyValleyField2012}
Mori, Masahiro, Karl MacDorman, and Norri Kageki. 2012. {``The {Uncanny
Valley} {[}{From} the {Field}{]}.''} \emph{IEEE Robotics \& Automation
Magazine} 19 (2): 98--100.
\url{https://doi.org/10.1109/MRA.2012.2192811}.

\bibitem[\citeproctext]{ref-muthukrishnanAmbiguityAversionPreference2009}
Muthukrishnan, A. V., Luc Wathieu, and Alison Jing Xu. 2009.
{``Ambiguity {Aversion} and the {Preference} for {Established
Brands}.''} \emph{Management Science} 55 (12): 1933--41.
\url{https://doi.org/10.1287/mnsc.1090.1087}.

\bibitem[\citeproctext]{ref-peracchioHowAmbiguousCropped1994}
Peracchio, Laura A., and Joan Meyers-Levy. 1994. {``How {Ambiguous
Cropped Objects} in {Ad Photos} Can {Affect Product Evaluations}.''}
\emph{Journal of Consumer Research} 21 (1): 190.
\url{https://doi.org/10.1086/209392}.

\bibitem[\citeproctext]{ref-porterAIgeneratedPoetryIndistinguishable2024a}
Porter, Brian, and Edouard Machery. 2024. {``{AI-generated} Poetry Is
Indistinguishable from Human-Written Poetry and Is Rated More
Favorably.''} \emph{Scientific Reports} 14 (1): 26133.
\url{https://doi.org/10.1038/s41598-024-76900-1}.

\bibitem[\citeproctext]{ref-shinUncannyValleyNo2019}
Shin, Mincheol, Se Jung Kim, and Frank Biocca. 2019. {``The Uncanny
Valley: {No} Need for Any Further Judgments When an Avatar Looks
Eerie.''} \emph{Computers in Human Behavior} 94 (May): 100--109.
\url{https://doi.org/10.1016/j.chb.2019.01.016}.

\bibitem[\citeproctext]{ref-taoExploringPersonaCharacteristics2022}
Tao, Yayi, Guangli Zhang, Di Zhang, Fang Wang, Yun Zhou, and Tao Xu.
2022. {``Exploring {Persona Characteristics} in {Learning}: {A Review
Study} of {Pedagogical Agents}.''} \emph{Procedia Computer Science} 201:
87--94. \url{https://doi.org/10.1016/j.procs.2022.03.014}.

\bibitem[\citeproctext]{ref-vukojicicImitationDrawingCan2023}
Vukojičić, Milić, Jovana Krstić, and Mladen Veinović. 2023. {``Imitation
{Drawing}: {Can We Spot} the {Difference Between Ai} and {Human
Generated Drawing}?''} In \emph{Proceedings of the {International
Scientific Conference} - {Sinteza} 2023}, 23--27. Beograd, Serbia:
Singidunum University.
\url{https://doi.org/10.15308/Sinteza-2023-23-27}.

\bibitem[\citeproctext]{ref-yiDeterminantsConsumerSatisfaction1993}
Yi, Youjae. 1993. {``The {Determinants} of {Consumer Satisfaction}: {The
Moderating Role} of {Ambiguity}.''} \emph{Advances in Consumer Research}
20: 502--8.

\end{CSLReferences}




\end{document}
